%! Characteristics of Polynomial Functions — Unit 3, Lesson 3.0 — Lesson Plan
\documentclass[10pt]{article}
\usepackage{algebra2-article}
\usepackage{algebra2-boxes}
\usepackage{graphicx}   % -article does NOT load graphicx; needed for thumbnails/screenshots
\graphicspath{{images/}}
\newcommand{\TallMath}[1]{$\displaystyle #1\rule[-1.4em]{0pt}{3.2em}$}
\newcommand{\CourseName}{Algebra 2: Shepherd}
\newcommand{\SchoolYear}{2026--2027}
\newcommand{\MeetingLength}{55 minutes}
\newcommand{\UnitNumberName}{Unit 3: Polynomial Functions \quad}
\newcommand{\LessonNumberName}{Lesson 3.0: Characteristics of Polynomial Functions}

\begin{document}
\begin{center}
    {\Large\bfseries \CourseName: \SchoolYear \\
    {\small\bfseries \UnitNumberName \LessonNumberName}}
\end{center}

\begin{tcolorbox}[colback=forestbg, colframe=forest, boxrule=0.9pt, arc=2mm,
    left=2mm, right=2mm, top=1.5mm, bottom=1.5mm]
    \textbf{Primary Objective:} Students can read the graph of a polynomial function --- a smooth
    curve that may rise and fall several times --- and name its characteristics: its \emph{degree}
    and \emph{leading coefficient} and the \emph{end behavior} they force, its \emph{turning
    points}, its \emph{relative (local)} and \emph{absolute (global)} maxima and minima, its
    \emph{domain}, \emph{range}, \emph{$y$-intercept}, and \emph{zeros} --- reading each zero's
    \emph{multiplicity} to decide whether the graph \emph{crosses} or \emph{touches} the $x$-axis
    --- the intervals where it is increasing/decreasing and positive/negative, and its
    \emph{symmetry} (even about the $y$-axis, \emph{odd} about the origin, or neither). This extends
    the parabola toolkit from Unit~2 to curves that turn more than once.
    \\[2pt]
    \textbf{Standards (2023 VA SOL):} \textbf{A2.F.2} --- investigate and analyze the
    characteristics of polynomial functions graphically. This Lesson~0 revisits \textbf{A2.F.2a}
    (domain, range, zeros, intercepts), \textbf{A2.F.2c} (increasing/decreasing intervals), and
    \textbf{A2.F.2d} (absolute max/min) on a polynomial, and \emph{introduces} \textbf{A2.F.2e}
    (location and value of \emph{relative} maxima/minima) and the \emph{degree $+$
    leading-coefficient} rule for \textbf{A2.F.2g} (end behavior). \textbf{A2.F.2b} (compare and
    contrast the characteristics of function families) frames the odd-degree vs.\ even-degree
    contrast.
\end{tcolorbox}

\renewcommand{\arraystretch}{1.4}
\begin{skillbox}[Priority Ideas \& Skills]{goldbox}
\begin{tabularx}{\linewidth}{@{}X|X@{}}
\textbf{Priority Ideas \& Skills} & \textbf{Key Understandings (the ``why'')} \\[2pt]
\hline
\vspace{-6pt}
\begin{itemize}[leftmargin=1.1em, nosep, after=\vspace{2pt}]
  \item Identify the \textbf{degree} and \textbf{leading coefficient} and state the \textbf{end behavior} from the degree--parity/sign rule.
  \item Count \textbf{turning points} and locate \textbf{relative} vs.\ \textbf{absolute} maxima/minima.
  \item Read \textbf{domain}, \textbf{range}, the \textbf{$y$-intercept}, and the \textbf{zeros}, using each zero's \textbf{multiplicity} to tell a \emph{cross} from a \emph{touch}.
  \item Give the \textbf{increasing/decreasing} and \textbf{positive/negative} intervals.
  \item Classify \textbf{symmetry}: \textbf{even} ($y$-axis), \textbf{odd} (origin), or neither.
\end{itemize}
&
\vspace{-6pt}
\begin{itemize}[leftmargin=1.1em, nosep, after=\vspace{2pt}]
  \item Only two facts --- \emph{is the degree even or odd?} and \emph{is the leading coefficient $+$ or $-$?} --- fix both end-behavior arms.
  \item A degree-$n$ polynomial turns \emph{at most} $n-1$ times; ``relative'' means locally highest/lowest, ``absolute'' means highest/lowest \emph{overall}.
  \item \textbf{Multiplicity} is written in the picture: \emph{odd} multiplicity crosses, \emph{even} multiplicity bounces (touches).
  \item ``Decreasing'' (direction) is \emph{not} ``negative'' (output sign) --- still true with several turns.
  \item Odd symmetry ($f(-x)=-f(x)$, a $180^\circ$ turn about the origin) is the new partner to even symmetry ($f(-x)=f(x)$).
\end{itemize}
\end{tabularx}
\end{skillbox}

\begin{skillbox}[{Vocabulary, Concepts, \& Theorems}]{sky}
\renewcommand{\arraystretch}{1.3}
\begin{tabularx}{\linewidth}{@{}l X@{}}
\textbf{Polynomial function} & $f(x)=a_nx^n+\dots+a_1x+a_0$ with whole-number exponents; a smooth, continuous curve (no breaks, holes, or sharp corners). \\
\textbf{Degree / leading coefficient} & The highest power $n$, and the coefficient $a_n$ of that term (written in standard form). \\
\textbf{End behavior} & What the two arms do as $x\to\pm\infty$; decided entirely by the parity of $n$ and the sign of $a_n$. \\
\textbf{Turning point} & A point where the curve changes from rising to falling or vice versa; a degree-$n$ graph has at most $n-1$. \\
\textbf{Relative (local) max/min} & The output at a turning point --- highest/lowest \emph{compared to nearby points}. \\
\textbf{Absolute (global) max/min} & The highest/lowest output of the \emph{whole} function; polynomials of odd degree have neither. \\
\textbf{Zero \& multiplicity} & An input where $f(x)=0$; its multiplicity is the exponent of that factor. Odd $\Rightarrow$ graph crosses; even $\Rightarrow$ graph touches (bounces). \\
\textbf{Even / odd symmetry} & Even: $f(-x)=f(x)$ (mirror across the $y$-axis). Odd: $f(-x)=-f(x)$ ($180^\circ$ rotation about the origin). \\
\end{tabularx}
\end{skillbox}

\begin{fixedskillbox}[Activate Prior Knowledge \& Spiral Review]{forestbg}
    The Warm-Up rehearses the skills this lesson stands on: \textbf{(1)} the sign of a power of a
    negative number (even power $\Rightarrow$ positive, odd power $\Rightarrow$ keeps the sign --- the
    seed of even/odd behavior and end behavior); \textbf{(2)} evaluating a cubic
    $f(x)=x^3-3x+2$ from a table (the outputs rise, fall, and rise --- \emph{two} turns, unlike a
    parabola); and \textbf{(3)} reading the \emph{zeros} straight off a \emph{factored} form. Debrief
    item~2 aloud: the table already shows more than one hill/valley, previewing turning points.
\end{fixedskillbox}

\begin{skillbox}[Hook]{forestbg}
    Show the graph of $g(x)=x^3-3x+2$. \textbf{``A parabola turns exactly once. How many times does
    this curve turn?''} (Twice --- a rise, a fall at the top of a hill, then a rise out of a valley.)
    Pose the questions that name every new feature: \textbf{``Where are the top of the hill and the
    bottom of the valley?''} (the \emph{relative} max $(-1,4)$ and \emph{relative} min $(1,0)$),
    \textbf{``At the far left and far right, where do the arms go?''} (down-left, up-right --- set by
    an \emph{odd} degree with a \emph{positive} lead), and \textbf{``Why does the curve bounce off the
    $x$-axis at $x=1$ but cut straight through at $x=-2$?''} (the factor $(x-1)^2$ has \emph{even}
    multiplicity; $(x+2)$ is \emph{odd}). Land the point: \emph{a polynomial's whole shape is encoded
    in its degree, its leading coefficient, and the multiplicities of its zeros.}
\end{skillbox}

\begin{skillbox}[Lesson]{forestbg}
\begin{multicols}{2}
\textbf{1. What makes a function a polynomial.} A \textbf{polynomial} is a sum of terms $a_kx^k$
with whole-number exponents; its graph is \emph{smooth} (no corners) and \emph{continuous} (no
breaks). In \emph{standard form} the \textbf{degree} is the highest power and the \textbf{leading
coefficient} is that term's coefficient. Contrast with non-examples ($x^{-1}$, $\sqrt{x}$) to fix
the ``whole-number exponents'' rule.

\textbf{2. End behavior --- the two-question rule.} The arms are decided by only two things.
\emph{Is the degree even or odd?} Even degree $\Rightarrow$ both arms go the \emph{same} way (like a
parabola); odd degree $\Rightarrow$ the arms go \emph{opposite} ways. \emph{Is the leading
coefficient $+$ or $-$?} Positive $\Rightarrow$ the right arm goes up; negative $\Rightarrow$ the
right arm goes down. Read all four cases off the parents $x^2,\,x^4$ (even) and $x^3,\,-x^3$ (odd).

\columnbreak

\textbf{3. Turning points \& extrema.} Between the arms a polynomial can rise and fall; each
change of direction is a \textbf{turning point}, and a degree-$n$ graph has \emph{at most} $n-1$ of
them. The output at a turning point is a \textbf{relative (local) max or min} --- highest/lowest
\emph{near there}. The \textbf{absolute (global)} max/min is the highest/lowest output of the
\emph{whole} function; stress that an odd-degree polynomial, whose arms split, has \emph{neither}.

\textbf{4. Zeros \& multiplicity.} On $g(x)=(x-1)^2(x+2)$ the zeros are $x=1$ and $x=-2$. The
\textbf{multiplicity} is the exponent of the factor: at $x=-2$ (multiplicity~1, \emph{odd}) the
graph \emph{crosses}; at $x=1$ (multiplicity~2, \emph{even}) it \emph{touches and turns back}
(bounces). Odd crosses, even bounces.

\textbf{5. Symmetry \& bringing it home.} Test symmetry with $f(-x)$: \textbf{even} if
$f(-x)=f(x)$ (mirror across the $y$-axis, only even powers, like $x^4$); \textbf{odd} if
$f(-x)=-f(x)$ ($180^\circ$ about the origin, only odd powers, like $x^3$). Re-read every feature of
$g$ and preview Lesson~3.1, \emph{operations with polynomials}.
\end{multicols}
\end{skillbox}

\begin{skillbox}[Active Monitoring]{redbox}
Circulate during the notes practice and Tier work. Watch for and correct:
\begin{itemize}[leftmargin=1.2em, nosep]
  \item calling a \emph{relative} max the \emph{absolute} max --- on an odd-degree graph the arm rises higher than any hill;
  \item reading multiplicity backward (thinking a \emph{cross} is the ``double'' root) --- odd crosses, even bounces;
  \item getting only one arm right --- for even degree both arms match, for odd degree they oppose;
  \item confusing ``decreasing'' with ``negative,'' now that the graph has several intervals of each;
  \item classifying a curve with both even and odd powers as ``odd'' --- odd symmetry needs \emph{only} odd powers.
\end{itemize}
Cold-call prompts: ``Even or odd degree --- and how do the arms tell you?'' ``Is that a relative or
an absolute minimum, and how do you know?'' ``Which zero makes the graph bounce, and why?''
\end{skillbox}

\begin{skillbox}[Group Work \& Differentiation]{redbox}
\textbf{Reading Polynomial Graphs} activity (tiered; mirrors the notes).
\begin{multicols}{3}
\textbf{Tier R --- Remediate.} From one pre-drawn cubic with three integer zeros, read the end
behavior, the number of turning points, the zeros, and the $y$-intercept, and mark each relative
max/min.

\columnbreak
\textbf{Tier A --- Approaching.} For a cubic and a quartic, give the full feature list --- degree
parity from the arms, turning points, relative/absolute extrema, zeros \emph{with cross/touch},
range, and symmetry --- and \emph{justify} the odd-symmetry classification with the $f(-x)$ test.

\columnbreak
\textbf{Tier E --- Extension.} A degree/leading-coefficient ``detective'': from a described arm
pattern and turning-point count, reason about the possible degree and sign; then interpret an
open-box volume model $V(x)=x(10-2x)(8-2x)$ in context (feasible domain, the meaning of the
max).
\end{multicols}
\end{skillbox}

\begin{skillbox}[Individual Work \& Assessment]{redbox}
\textbf{Exit Ticket} (independent, no notes): from one graphed cubic $k(x)=(x+1)^2(x-2)$, state the
end behavior, the turning points and relative extrema, the zeros \emph{with their cross/touch
behavior}, and the $y$-intercept; \textbf{classify} the symmetry with justification; and answer an
\textbf{SOL-style multiple-choice} item separating end behavior, multiplicity, and relative vs.\
absolute extrema. Collect and sort into ``got it / multiplicity slip / relative-vs-absolute slip /
end-behavior slip'' to decide whether Lesson~3.1 opens with a quick re-teach.
\end{skillbox}

\begin{skillbox}[Reinforcement \& Extension]{goldbox}
\textbf{Homework:} from a factored form, list the zeros with multiplicity and the cross/touch
behavior and state the end behavior; read the full feature list from a cubic and from an even
quartic; and one justification item on why an odd-degree polynomial has no absolute maximum.
\textbf{Extension:} interpret an open-box volume model in context, including its feasible domain.
\textbf{Preview:} Lesson~3.1, \emph{Operations with Polynomials} --- adding, subtracting, and
multiplying the expressions whose graphs we just learned to read.
\end{skillbox}
\end{document}
